% sec_representation.tex — the representation theorem
\section{Aggregation: the Representation Theorem}\label{sec:representation}

The composite indicator has the gated form
\begin{equation}\label{eq:composite}
A \;=\; \Avia\cdot F_\varsigma(\Aorg,\Agoal,\Aexec,\Aepi,\Aref),\qquad
F_\varsigma(x) = \Bigl(\textstyle\sum_i w_i\,x_i^{(\varsigma-1)/\varsigma}\Bigr)^{\frac{\varsigma}{\varsigma-1}},
\end{equation}
where the CES family $F_\varsigma$ is parameterized by the elasticity
$\varsigma$ (equivalently the power-mean exponent $p=(\varsigma-1)/\varsigma$;
the elasticity symbol is $\varsigma$ throughout, to be distinguished from the
scaffolding intervention class $\sigma$). This section
shows that the main case $\varsigma=1$---the weighted geometric mean
$F(x)=\prod_i x_i^{w_i}$---is \emph{characterized} by the axioms, twice over:
forced relative to the framework's stated normative commitments (which are
defended and priced below and in Section~\ref{sec:objections}), not by
measurement theory alone.

\subsection{Statement}

\begin{definition}[(A1)--(A6)]\label{def:a16}
For $F:(0,1]^n\to(0,1]$, continuously extended to $[0,1]^n$:
\textbf{(A1)} continuity (AX4, regularity); \textbf{(A2)} strict increase in
each coordinate (AX4); \textbf{(A3)} \emph{biseparability}: for every
non-trivial partition
$\{K,L\}$ of the coordinates there are continuous increasing $G_{KL}$ and
sub-aggregations $F_K,F_L$ of the same kind with
$F=G_{KL}(F_K(x_K),F_L(x_L))$ (AX7)---the precise form of ``the order of
subgroup aggregation does not matter''; \textbf{(A4)} zero annihilation (AX5);
\textbf{(A5)} homogeneity of degree one, $F(tx)=tF(x)$ whenever
$tx\in(0,1]^n$ (AX8; a local property on the bounded cube, extended to
$\mathbb{R}_{++}^n$ by Lemma~\ref{lem:homog-ext}); \textbf{(A6)}
coordinatewise scale invariance of the induced order (AX6b). Each condition
is thus an axiom of Section~\ref{sec:axioms} with a stated normative defense;
none is an undefended technical convenience. The minimal sufficient axiom
set, however, is the invariance trio $\{$AX4, AX6b, AX8$\}$: route B below
uses neither (A3) nor (A4), and the product form it returns satisfies both
(Remark~\ref{rem:routeb-minimal}).
\end{definition}

\begin{theorem}[Representation theorem]\label{thm:rep}
If $F$ satisfies \textup{(A1)--(A6)}, then $F(x)=\prod_i x_i^{w_i}$ with
$w_i>0$, $\sum_i w_i=1$, uniquely determined by $F$ via the elasticities
$w_i=\partial\log F/\partial\log x_i$.
\end{theorem}

\begin{proof}[Proof sketch (two routes, independent toolchains)]
\emph{Route A} (aggregation theory): (A1)--(A3) give a quasilinear
representation $F=\phi^{-1}(\sum_i w_i\phi(x_i))$ by the
Nagumo--Kolmogorov--de Finetti theorem ($n=2$) and the
biseparable/separability theory of \citet{debreu1960} with the
reduction-to-common-generator step of Acz\'el
\citep{aczel1966,aczel_dhombres1989,gorman1968}; (A4) forces
$\phi(0^+)=-\infty$; (A5) solves an Acz\'el--Kolmogorov functional equation
self-containedly (weighted Jensen $\Rightarrow$ affine; multiplicative Cauchy
$\Rightarrow$ $\Psi(x)=x^\rho$), leaving the power-mean family $M_p$ with
$p\le0$; (A6) excludes $p<0$ by an explicit order-reversal witness, leaving
$p=0$. \emph{Route B} (scale-invariance route, no separability): on the
bounded cube (A5) is a local property, so we first extend $F$ homogeneously
to $\mathbb{R}_{++}^n$ (Lemma~\ref{lem:homog-ext}); in
log-coordinates $g(u)=\log F(e^u)$, (A6) becomes translation invariance of the
induced order; the set $D$ of order-preserving shifts is a closed convex cone
with $D\cup(-D)=\mathbb{R}^n$; a half-space lemma on its quotient by the
indifference subspace forces that quotient to be one-dimensional, so
$g(u)=w\cdot u$; homogeneity and $F(1)=1$ then give $\sum_iw_i=1$ and the
product form. Full proofs: Appendix~\ref{app:representation}.
Because route B does not use (A3) or (A4), the theorem's minimal sufficient
set is the invariance trio (AX4, AX6b, AX8), and biseparability and zero
annihilation are consequences rather than independent commitments
(Remark~\ref{rem:routeb-minimal}).
\end{proof}

\paragraph{The function form is derived; the weights are chosen.}
Theorem~\ref{thm:rep} pins down the \emph{form} of $F$; the elasticities
$w_i=\partial\log F/\partial\log x_i$ are then an identity of the resulting
composite, not an estimation procedure---regressing $\partial\log
F/\partial\log x_i$ presupposes $F$, hence the $w_i$ (a circularity we do not
commit). Two normalization details are explicit: $F(1,\dots,1)=1$ is what
fixes $\sum_iw_i=1$ (route A inherits the Nagumo--Kolmogorov normalization;
route B assumes it), and without that normalization the product carries a
multiplicative constant, $c\prod_ix_i^{w_i}$ with $c>0$. Three non-circular determination protocols are compatible with the
framework: (a) \emph{equal weights} as the transparent default, defended
normatively---each channel is a necessary condition of autonomy, and nothing
in the framework licenses ex ante inequality (the benchmark uses this);
(b) \emph{independently elicited} importance weights (expert or stakeholder
annotation), registered as application-layer input; (c) \emph{behaviorally
measured} weights, recovered from the agent's own trade-offs under controlled
single-channel perturbations. Whichever protocol is used, reporting a
weight-sensitivity band is an obligation: Section~\ref{sec:experiments}
recomputes every campaign verdict under five alternative weight vectors
(ordering Kendall $\tau\ge0.93$ against equal weights; zero verdict flips).

That two derivations---axiomatic aggregation theory and measurement scale
theory, sharing the invariance axioms but independent in mathematical
toolchain---converge on the same product form is our answer to
the standing objection that composite aggregation is arbitrary
(\citealp{oecd2008}; Section~\ref{sec:objections}).

\paragraph{Attribution and the locus of novelty.}
The scale-invariance characterization is classical, and we locate our
additions exactly. \citet{ebert_welsch2004} proved, in a social-choice
analysis of environmental indices, that meaningful aggregation of ratio-scale
indices under scale-invariance requirements yields geometric
(Cobb--Douglas-type) forms---the result the OECD handbook inherits as its
scale-type discipline \citep{oecd2008}. Route B is best read as the
bounded-cube, order-invariance form of that characterization, re-derived
with the audit setting's quantifiers (invariance of the \emph{induced order}
under coordinatewise rescaling, localized to $(0,1]^n$ by the homogeneous
extension of Lemma~\ref{lem:homog-ext}). What this paper adds is not the
product form itself but: (i) the \emph{minimal-trio} analysis---zero
annihilation and biseparability are consequences of $\{$AX4, AX6b, AX8$\}$,
not independent commitments (Remark~\ref{rem:routeb-minimal}); (ii) the
harmonic counterexample isolating exactly which commitment excludes the
harmonic family (Proposition~\ref{prop:harmonic}); (iii) the convergence of
two independent toolchains on one form; and (iv) the audit embedding---the
viability gate as an explicit departure from the theorem's hypotheses, the
non-circular weight-determination protocols, the falsification windows, and
the mandatory axiom-relaxation bands of Section~\ref{sec:experiments}.
Novelty claims in this paper are to be read with this scope.

\begin{proposition}[The harmonic counterexample]\label{prop:harmonic}
The harmonic mean $H(x)=(\sum_i w_i x_i^{-1})^{-1}$ satisfies continuity,
strict monotonicity, zero-annihilation, biseparability (quasi-arithmetic
generator $-1/x$), and homogeneity, but violates \textup{(A6)}: with
$n=2$, $w=(\tfrac12,\tfrac12)$, $x=(0.25,1)$, $y=(0.5,0.5)$,
$\lambda=(1,0.25)$, one has $H(y)=0.5>H(x)=0.4$ yet
$H(\lambda\odot y)=0.2<H(\lambda\odot x)=0.25$. Hence the invariance axiom
is indispensable: the remaining five conditions do not characterize the
geometric mean, and (A6) is exactly what excludes the harmonic family.
\end{proposition}

\begin{corollary}[The admissible elasticity is a singleton]\label{cor:sigma-band}
$p\le0\iff\varsigma\in(0,1]$ is the domain admitted by
\textup{(A1)}--\textup{(A5)} alone; \textup{(A6)} then excludes every $p<0$
(Lemma~\ref{lem:ax6b}), so the full axiom set admits exactly
$\varsigma=1$---the weighted geometric mean. $\varsigma\to0$ (Leontief/min) is
excluded by \textup{(A2)}; $\varsigma\to\infty$ (arithmetic) by
\textup{(A4)}; $\varsigma>1$ violates zero-annihilation. The framework's
reporting discipline follows from this, with the epistemic status made
explicit: all conclusions are computed at the admissible point $\varsigma=1$,
and the bands $\varsigma\in[0.5,1)$ (relaxing \textup{(A6)}) and
$\varsigma=2$ (relaxing \textup{(A4)}/AX5) are reported as
\emph{axiom-relaxation robustness checks} of exactly equal standing
(Section~\ref{sec:experiments})---they probe whether verdicts survive in a
neighborhood of the forced form; they are not axiom-admissible alternatives.
\end{corollary}

\paragraph{The normative commitments, and what rejecting one costs.}
The theorem is conditional on the invariance trio; the following map records
what each rejection buys and what it destroys, so that the normative choice is
made with its price visible.
\begin{center}\small
\begin{tabularx}{\textwidth}{@{}l >{\raggedright\arraybackslash}X
>{\raggedright\arraybackslash}X@{}}
\toprule
Rejected commitment & Admissible family & What is lost \\
\midrule
AX6b (scale invariance) & power means $M_p$, $p\le0$ (harmonic included) &
unique form; order can flip under unit rescaling (Prop.~\ref{prop:harmonic});
cross-calibration comparisons ill-defined; T4/T5/T7 conclusions become
family-relative \\
AX8 (homogeneity) & level-dependent weightings (AX4--AX7 kept) &
composite scale no longer tracks channel scales; T4's log-linear LP and T7's
$(1-m)^{w_{org}}$ law lose their exact form; T6's product floor becomes
level-dependent \\
AX4 (strict monotonicity) & sub-support products, $\min$, discontinuous rules &
strict improvement fails; T1b positivity fails; coverage/T3 strict-decrease
conclusions weaken to non-strict; T5's ``iff'' becomes one direction \\
AX5 (zero annihilation) & not independently rejectable: within the invariance
class AX5 $\iff$ strict monotonicity (Remark~\ref{rem:routeb-minimal}); a
non-annihilating stance (e.g.\ arithmetic aggregation) must also give up AX6b &
the conjunctive-necessity semantics; T5's zero-disjunction reading; the
$\varsigma=2$ stress arm's failure mode (complete capture reads
$A=0.1652>0$) \\
\bottomrule
\end{tabularx}
\end{center}
The AX5 stance---any channel at zero annihilates autonomy---is the strongest
of the aggregation commitments, and it is defended as conjunctive necessity:
each channel is a necessary condition of the operational definition, so no
surplus on one channel compensates a zero on another. The
multidimensional-poverty literature supplies the precedent
\citep{alkire2011}, and the benchmark supplies the sensitivity obligation:
the $\varsigma=2$ arm exhibits exactly what zero annihilation buys
(complete capture at $A=0$ versus $A=0.1652$ under the relaxed form). The
cost is registered: the framework does not permit compensation for a dead
channel; applications that require it must move to the relaxation band and
report the move. Everything not listed in the third column survives the
corresponding rejection: in particular, T6's comparison-principle margins and
the identifiability analysis do not depend on the aggregation form, and T2's
attribution cover survives for any strictly monotone $F$ (the strict-decrease
readings weaken only if AX4 itself is dropped).

\paragraph{Why these two invariance commitments at all (the positive defense).}
The rejection map prices refusal; here we state what acceptance buys, since
AX6b and AX8 carry the exclusion load and deserve affirmative defenses rather
than silent uptake. \emph{AX6b (coordinatewise scale invariance).} Under
AX6a's canonical normalizers, coordinate rescaling is not an admissible
transformation in the strict scale-type sense---we registered this in
Section~\ref{sec:axioms} and do not claim measurement-theoretic necessity.
The commitment is instead \emph{defense in depth against the audit's
participants}: normalizers are themselves protocol choices made by fallible
instrument designers, and AX6b guarantees that no unilateral
re-instrumentation of a single channel---a vendor recalibrating the goal
metric, a platform renormalizing its empowerment probe---can flip the
ordering of two audited agents. That is the operational content of ``no party
to the audit can change the verdict by changing a ruler''; it is also what
makes cross-deployment comparisons meaningful when channels were normalized
by different teams. The price is visible in the map: the harmonic family,
which some welfare traditions prefer for its penalty on imbalance, is
excluded (Proposition~\ref{prop:harmonic}). \emph{AX8 (homogeneity).} Audit
language is percentage language: governance asks ``how much autonomy did the
interaction cost?'' and expects the answer to compose---equal proportional
damage to every channel should mean equal proportional damage to autonomy.
Homogeneity is exactly what ties the composite's scale to the channels'
ratio scales and makes marginal rates of substitution between channels
level-independent; without it, a ``$20\%$ autonomy loss'' has no
channel-consistent meaning. The price is equally visible: level-dependent
weighting schemes---weights that shift with the autonomy level---are
excluded, and applications that want them must move to the relaxation band
and report the move (Corollary~\ref{cor:sigma-band}). Both defenses are
offered as arguments, not proofs; the sensitivity obligations of
Section~\ref{sec:experiments} exist precisely because they may be rejected.

\subsection{Gate structure and dependence}

\begin{assumption}[Via-Ind: conditional independence of channel degradations]
\label{ass:via-ind}
Given viability, channel degradations are conditionally independent.
\end{assumption}

Writing $\Avia$ as an external multiplicative gate corresponds to
(Via-Ind). The gate's axiom status is
one-sided and we state it exactly: as an $(n{+}1)$-th geometric factor with
external weight $1$ outside the $\sum_iw_i=1$ normalization, the gated form
$\tilde F(x_{\mathrm{via}},x)=x_{\mathrm{via}}\cdot\prod_ix_i^{w_i}$ preserves
continuity, strict monotonicity, biseparability, zero-annihilation, and
coordinatewise scale invariance, but it is homogeneous of degree
$1+\sum_iw_i=2$ in the extended vector, so \emph{homogeneity \textup{(A5)} is
deliberately not preserved}: the gate is an explicit departure from the
hypotheses of Theorem~\ref{thm:rep}, defended by its semantics (a viability
cut scales the whole once, independently of any proportional rescaling of the
five channels) rather than absorbed into the theorem. \emph{Scope statement}:
Theorem~\ref{thm:rep} applies to the five-channel aggregator $F_\varsigma$;
the full index is the composite design
$A=\Avia\cdot F_\varsigma$, and every ``forced'' claim in this paper is to be
read with that scope.
When (Via-Ind) fails---an attack striking channel clusters at once---nested
CES aggregation is automatically compatible with the axioms because (A3) is
exactly subgroup-order independence; and with marginals known but dependence
unknown, $\mathbb{E}[F]$ admits computable Fr\'echet/comonotonicity bounds
(Appendix~\ref{app:copula}). Component uncertainty biases the composite
\emph{downward} in expectation (Jensen gap, Lemma~\ref{lem:jensen-gap}), so
black-box interval mode that substitutes expected components overestimates $A$
by a signed amount.
